\documentclass[12pt]{article}

\usepackage{graphicx}

\usepackage{makecell}
\usepackage[figuresright]{rotating}
\usepackage{multicol, caption}
\usepackage{authblk}
\usepackage[margin=1.0in]{geometry}
\usepackage[export]{adjustbox}

\textheight=210mm
\textwidth=155mm

\hoffset=-1mm
\voffset=-10mm

\begin{document}
Submission MS \#CHA19-AR-00322


\begin{center}
{\bf Response to comments of reviewer \#3}\\
\vspace{2mm}
Concerning paper ``Generalised early warning
signals in multivariate and gridded data with an application to
tropical cyclones''\\[5pt]
by {\it J.~Prettyman, T.~Kuna, V.~N.~Livina}
\end{center}

\vspace{3mm}

We thank the reviewer for the time taken to read through this submission and provide positive constructive notes for improvement. We will reply to the comments here in the order in which they are given.

\begin{enumerate}
\item \textit{I have more reservations regarding the approach taken from [22] -- please consider dropping or shortening this part, since indeed the relation of Jacobian eigenvalues with the motion of a cyclone seems a bit far fetched.}
	
\textbf{reply (1):} We have now added a paragraph to the manuscript (bottom of page 4) explaining our decision to test existing tipping point analysis techniques in a physical system which exhibits a tipping point with no evidence of genuine bifurcational behaviour. In particular the generality of the principles of critical slowing down and increase in autocorrelation, which is also mentioned (page 5, right column) in the context of the Jacobian eigenvalues method. 


\item \textit{I am not sure if the references to dynamical systems' examples and their bifurcations in Sections III and IV are really so helpful for this paper. Maybe (parts of) that material could be moved into the appendix / supplements, so that the focus is more on data analysis and modelling.}
	
\textbf{reply (2):}	We appreciate this comment from a different perspective, but we believe the result in section III is a key novel result in the paper. To our knowledge, the effect on autocorrelation in a bifurcating system of the EOF technique has not before been studied and so, before using this technique in our later analysis of the sea-level pressure data, we believed it important to include this working in full. Section IV was expanded at the request of another reviewer, and we agreed at the time that an example of the method at work would be a useful addition. We would prefer, therefore, to leave this decision to the editor. Again, we might hope that the extra explanation in section V of the applicability of this previous research [22] to tipping points \emph{in general} makes the inclusion of section IV more acceptable. 
	
\item \textit{Minor points: \begin{itemize}
		\item Please use "SLP" in Fig. 1 and its caption consistently (not sometimes "slp"). 
		\item Please check for typos, e.g., the beginning of Section IV: "in Ref. 22, a method is introduce to anticipate ..." and on page 7: "We have used values such that a =0.4 for t The part of the noise signal ..." 
	\end{itemize}
}

\textbf{reply (3):}	We have now read through the manuscript more closely and fixed a few minor errors. Thank you in particular for drawing our attention to the second point here: a significant piece of text (present in previous versions) had gone missing at this point in the manuscript (now included on page 7 - right column). Hopefully this section now makes sense. 

\end{enumerate}

We thank the reviewer again for their time in responding carefully, and positively, to this submission. We are happy that the reviewer has deemed the paper suitable for publication in the Chaos magazine and hope that the minor adjustments made to the manuscript are considered improvements.

\vspace{5mm}

\noindent
Yours faithfully,

\vspace{1mm}

\noindent
\it J.~Prettyman, T.~Kuna, V.~N.~Livina

\end{document}
